% !TeX root = ../sections/02_exercises.tex

\begin{exercise}
	R-7. $A=\mathbb{Z}\oplus\mathbb{Z}\oplus\mathbb{Z}$ を階数 $3$ の
	$\mathbb{Z}$-自由加群とし，その部分加群 $B,C$ を次のように置く：
	\[
	B:=\mathbb{Z}(1,1,0)+\mathbb{Z}(0,1,1)+\mathbb{Z}(1,0,1),
	\]
	\[
	C:=\mathbb{Z}(1,1,0)+\mathbb{Z}(0,1,1)+\mathbb{Z}(2,5,3).
	\]
	\begin{enumerate}
		\item $B$ と $C$ の階数を求めよ。
		\item $A/B$ と $A/C$ の不変系を求めよ。
	\end{enumerate}
\end{exercise}

\begin{background}
	$\mathbb{Z}^3$ の部分加群が複数のベクトルで生成されているとき，
	それらを列ベクトルとして並べた整数行列を考える。
	
	部分加群の階数は，その行列の $\mathbb{Q}$ 上での階数に等しい。
	また，商加群の不変系は Smith 標準形によって求められる。
	
	特に，$3\times 3$ 行列の行列式が $0$ でなければ，
	その列ベクトルは $\mathbb{Q}$ 上一次独立であり，部分加群の階数は $3$ である。
	一方，行列式が $0$ の場合は，列ベクトル間の一次関係を調べる必要がある。
\end{background}

\begin{strategy}
	$B$ については，生成ベクトルを列に並べて
	\[
	M_B=
	\begin{pmatrix}
		1 & 0 & 1\\
		1 & 1 & 0\\
		0 & 1 & 1
	\end{pmatrix}
	\]
	を考える。
	この行列式は
	\[
	\det M_B=2
	\]
	であるから，$B$ の階数は $3$ である。
	
	さらに，$M_B$ の成分には $1$ が含まれ，$2$ 次小行列式にも $1$ が現れるため，
	Smith 標準形は
	\[
	\operatorname{diag}(1,1,2)
	\]
	である。
	したがって
	\[
	A/B\cong \mathbb{Z}/2\mathbb{Z}
	\]
	となる。
	
	$C$ については
	\[
	M_C=
	\begin{pmatrix}
		1 & 0 & 2\\
		1 & 1 & 5\\
		0 & 1 & 3
	\end{pmatrix}
	\]
	を考える。
	第 $3$ 列は
	\[
	(2,5,3)=2(1,1,0)+3(0,1,1)
	\]
	と表される。
	したがって，$C$ は実質的に
	\[
	(1,1,0),\quad (0,1,1)
	\]
	で生成され，階数は $2$ である。
	
	この $2$ 本の生成ベクトルを列に並べた行列には，
	$2$ 次小行列式として $1$ が現れる。
	よって，この部分加群は $\mathbb{Z}^3$ の中で原始的であり，
	商加群にはねじれ部分がない。
	
	したがって
	\[
	A/C\cong \mathbb{Z}
	\]
	である。
\end{strategy}